\documentclass[11pt,a4paper]{article}
\usepackage[utf8]{inputenc}
\usepackage[T1]{fontenc}
\usepackage[french]{babel}
\usepackage{amsmath,amssymb}
\usepackage{geometry}
\geometry{margin=2cm}

\begin{document}

\begin{center}
    {\Large \textbf{Classe de Seconde -- Mathématiques}} \\[0.3em]
    {\large \textbf{Exercice : Fonctions et variations}}
\end{center}

\vspace{0.5cm}

Soit $f$ la fonction définie sur $\mathbb{R}$ par $f(x) = 2x^2 - 8x + 6$.

\begin{enumerate}
    \item \textbf{Forme canonique :}\\
    Démontrer que pour tout réel $x$, on a : 
    $$f(x) = 2(x - 2)^2 - 2$$

    \item \textbf{Sommet et extremum :}\\
    En déduire les coordonnées du sommet de la parabole représentant $f$ ainsi que son extremum sur $\mathbb{R}$.

    \item \textbf{Tableau de variations :}\\
    Préciser les intervalles de croissance et de décroissance de la fonction $f$.

    \item \textbf{Résolution d'équation :}\\
    Résoudre l'équation $f(x) = 6$.
\end{enumerate}

\vspace{0.5cm}
\hrule
\vspace{0.5cm}

\begin{center}
    {\large \textbf{Corrigé détaillé}}
\end{center}

\begin{enumerate}
    \item \textbf{Forme canonique}\\
    Développons l'expression proposée :
    $$2(x - 2)^2 - 2 = 2(x^2 - 4x + 4) - 2 = 2x^2 - 8x + 8 - 2 = 2x^2 - 8x + 6 = f(x)$$
    L'égalité est ainsi démontrée pour tout réel $x$.

    \item \textbf{Sommet et extremum}\\
    La forme canonique est de la forme $a(x - \alpha)^2 + \beta$ avec $a = 2$, $\alpha = 2$ et $\beta = -2$.
    Le coefficient $a = 2 > 0$, donc la parabole est tournée vers le haut.
    Le sommet $S$ a pour coordonnées $(2 ; -2)$ et la fonction admet un minimum égal à $-2$ atteint en $x = 2$.

    \item \textbf{Tableau de variations}\\
    La fonction est strictement décroissante sur $]-\infty ; 2]$ et strictement croissante sur $[2 ; +\infty[$.

    \item \textbf{Résolution de $f(x) = 6$}\\
    En utilisant l'expression de départ $2x^2 - 8x + 6 = 6$, on obtient $2x^2 - 8x = 0$, soit $2x(x - 4) = 0$.
    Les solutions sont $x = 0$ et $x = 4$. L'ensemble des solutions est $\{0 ; 4\}$.
\end{enumerate}

\end{document}
